% 第 4 章　方向也能计算：向量
\chapter{方向也能计算：向量}

\step{反常识开场}

\begin{mainline}
一个没有风的雨天，你站在路边看雨：雨丝几乎竖直地落下来，伞直直地撑在头顶就够用。然后你坐进一辆行驶的汽车，看侧面的车窗——玻璃上的水痕不再竖直，而是斜斜地从前上方向后下方拖。车开得越快，水痕越平，几乎横着走。请你停下来想一想这个画面有多奇怪：雨明明是竖直落的，没有任何人从侧面推它，是谁把雨“掰斜”的？

第二个谜题发生在人行道上。打开手机的计步和地图，先朝正东走三百米，再左转朝正北走四百米。计步器忠实地报告：您走了七百米。可地图上从起点到终点的直线距离只有五百米。你的确一步一步走完了七百米，地图的测量也没错——那“消失”的两百米去哪了？

第三个谜题只需要两只弹簧秤。用它们同时水平拉一个小铁环：两只秤各读五十牛，方向相同时，环受到的“总拉力”是一百牛；方向相反时，总拉力是零；而让两个拉力夹成六十度角时，总拉力约是八十七牛——不多也不少。五十加五十，怎么会是八十七？弹簧秤又没坏。

三个谜题表面互不相干，背后却是同一件事：我们从小学习的算术只管“多少”，不管“朝哪”。而雨的下落、人的行走、秤的拉力，全都带着方向。方向不是可以随手扔掉、算完再贴回去的包装纸——它本身就是数值的一部分。这一章要发明的，是一门把方向也算进去的算术。学会了它，你不但能解开这三个谜题，还会明白帆船为什么能比风跑得快、飞机为什么常常斜着机头飞直线，以及望远镜为什么要像车窗一样“斜着看”星星。
\end{mainline}

\step{先猜一猜}

在往下读之前，请像前几章一样，把直觉押在桌面上。猜错不丢人——本章要修理的正是直觉。

第一题。回到雨天的侧车窗：水痕斜向后下方，意思是水痕的上端靠车头、下端靠车尾，还是反过来？凭感觉画一下，不用急着讲道理。

第二题。一只小船要横渡一条河，水流稳定地向下游冲。船夫把船头\emph{正对}着对岸划（船相对水的运动方向与河岸垂直），船会在正对岸登陆，还是在下游登陆？如果想正好在对岸登陆，船头该朝上游偏，还是朝下游偏？

第三题。你拖着带拉杆的行李箱赶路，拉杆与地面夹着一个角。在拉力大小不变的前提下，拉杆抬得越陡，箱子向前的“劲头”是越大、越小，还是不变？直觉多半会说“使的劲一样大，效果就一样”——请先记住这句话，本章要当面拆穿它。

写下你的三个答案。等到本章结尾，我们逐一对账。

\step{动手实验}

\begin{experiment}[影子与拼路：方向的两张面孔]
\textbf{材料}：手机的手电筒功能、一支铅笔、一张白纸、一把尺子；一段能走开的空地（客厅地砖就很好用）、一根两三米长的绳子。

\textbf{步骤一（影子实验）}：把手机平放在桌沿，打开手电，让光水平地射向竖在桌面上的白纸。把铅笔竖在光路与纸之间，先让铅笔与光线\emph{垂直}（横躺着挡住光），在纸上描下影子的长度；然后让铅笔慢慢转向与光线\emph{平行}的方向。你会看到影子从“一整支铅笔长”逐渐缩短，最后缩成一个点。铅笔全程没有变形，变的是它在光线方向上的“存在感”。用尺子量两三种角度下的影长，记下来。

\textbf{步骤二（拼路实验）}：在空地上用脚步拼一个直角：从起点朝一个方向走三大步，左转九十度，再走四大步。把绳子从起点直接拉到终点，沿着你的脚步量一量——绳长大约是五大步。三大步加四大步，“直线回去”却只要五大步。再走一次，这次第二步只走两大步，绳长会变成多少？先猜再量。

\textbf{记录}：两个实验各记一个问题。影子里：铅笔的长度没有变，为什么影子会变？拼路里：脚步总数没有变，为什么“回到起点的直线距离”不等于脚步总数？把问题写下来，本章会逐个回答。这两个实验其实是一对：一个讲“一个带方向的量在指定方向上显露出多少”，另一个讲“两段带方向的路怎样合成一段”——它们正是本章两大工具的雏形。
\end{experiment}

\step{旧模型遇到麻烦}

先给旧模型画个像。我们从小学算术里自然长出一套世界观：凡是能测量的东西，归根结底都是一个数；测量就是读数，计算就是把数加减乘除。温度是数，质量是数，钱包里的余额是数。这套模型管用的场合多得惊人：两杯五十摄氏度的水兑在一起不会变成一百度——温度不能乱加；但两千克的米再加三千克，一定是五千克，从不失手。

可是麻烦来了，而且一来就是一串。

麻烦一：拼路实验。向东三步加向北四步，结果不是七步，而是沿斜线的五步。“三步加四步”这件事本身就有两种答法，取决于你问的是“总共走了多少路”，还是“现在离家多远”。旧模型里根本没有这两个问题的区别——它只有一个“总数”。

麻烦二：两只弹簧秤。五十牛加五十牛，答案在一百牛与零之间随角度任意变化。如果说弹簧秤读的是“数”，这些数公然违抗加法。旧模型在这里彻底破产：它连“总拉力是多少”这个最基本的问题都给不出确定的算法。

麻烦三：开场的雨。竖直下落的雨，没有任何水平方向的来路，却在车窗上留下水平方向的分量。旧模型连“这个分量从哪来”都无法表述，因为在它的语言里，运动只有快慢，没有朝向。

三条麻烦指向同一个诊断：\emph{有一类量，方向就是它数值的一部分}。把方向扔掉，剩下的数自然算不对账。这不是说旧算术错了——米还是五千克，账还是那些账——而是说它只管得住世界的一半。历史上，水手和商人早就熟练地使用方向（罗盘导航比微积分古老得多），但把方向本身变成可以做加减乘除的对象，是相当晚近的事：十八世纪末，测量员韦塞尔为了算地图，第一次把平面上的“带方向的线段”当成可以运算的数；十九世纪中叶，哈密顿为了描述三维转动发明了四元数，吉布斯和赫维赛德随后从中剥离出今天教科书里的向量算术。也就是说，向量不是从天上掉下来的“显然工具”，而是一批人被地图、转动和电磁学逼出来的发明——麦克斯韦整理电磁方程组时，正是这套新语言帮他把几十个分量方程压缩成干净的几行。物理学家几乎是立刻发现，他们手里最难缠的那些量——位移、速度、力——全都天生适合这门新算术。

\step{发明新概念}

\begin{mainline}
\textbf{第一个新概念：标量与向量。}把测量对象分成两类。只回答“多少”的叫\emph{标量}：温度、质量、路程、话费，一个数（带单位）说完。必须同时回答“多少”和“朝哪”的叫\emph{向量}：位移（从哪里到哪里）、速度（多快、朝哪）、加速度、力。判断一个量是不是向量，不看它高不高级，只看一个方向反转的测试：把它朝反方向来一份，效果会不会改变？五千克米反过来还是五千克米，标量；五十牛的推力换成反方向，门从被推开变成被拽上，向量。

\textbf{第二个新概念：位移不是路程。}这是本章最重要的一次分家。路程是你脚下实际碾过的总长度，标量；位移是从起点指向终点的直线段，带方向，向量。拼路实验里，七步是路程，沿斜线的五步是位移。“消失的两百米”从来没有消失——“走了多少路”和“离开了多远”本来就是两道不同的题，只是它们恰好在走直线时答案相同，骗我们以为是一回事。

\textbf{第三个新概念：向量是一支“会平移的箭头”。}数学上，向量用一支箭头表示：长度代表大小，朝向代表方向，仅此而已。它像什么？像一句导航指令——“向东三百米”。这句话谁拿到、在哪个路口执行，效果都一样；指令本身不含起点。它不像什么？它不像地图上的一个点（向量的起点不在向量里，北京城东三百米的箭头和上海城东三百米的箭头是\emph{同一个}向量，画在哪儿都行）；也不像一根钉死在原地的木桩（平行挪动箭头，箭头还是原来那个向量）。正因为向量不依附于位置，我们才能把作用在同一物体上的几个力随手搬到一起比较、拼接。

\textbf{第四个新概念：向量加法——首尾相接。}向东走三百米、再向北走四百米，合起来的效果是什么？把第二支箭头的尾巴接到第一支箭头的头上，从第一支箭头的尾巴指向第二支箭头的头，就是答案：位移 \(\vec{A}\) 接着位移 \(\vec{B}\)，合位移写作 \(\vec{A}+\vec{B}\)。这叫三角形法则；把两支箭头从同一起点画出来、补成平行四边形取对角线，是同一件事的另一种画法。请注意这里的分工，本书会一直坚守它：\emph{“箭头首尾相接”是数学定义}，而“位移按这个规则合成”是几何上可以直接验证的事实；“力也按这个规则合成”则是\emph{实验事实}——两只弹簧秤拉环、再换一只秤单独拉到同样效果，读数对比，次次如此。定义是我们发明的，事实是世界投票的。
\end{mainline}

\begin{center}
\begin{tikzpicture}[scale=0.9]
  \coordinate (O) at (0,0);
  \coordinate (E) at (3,0);
  \coordinate (N) at (3,4);
  \draw[-{Stealth[length=3mm]},very thick,blue!70!black] (O) -- (E) node[midway,below=2pt] {向东 \(3\) 步};
  \draw[-{Stealth[length=3mm]},very thick,blue!70!black] (E) -- (N) node[midway,right=2pt] {向北 \(4\) 步};
  \draw[-{Stealth[length=3mm]},very thick,red!70!black] (O) -- (N) node[midway,above left=-1pt] {合位移 \(5\) 步};
  \fill (O) circle (1.5pt) node[below left=1pt] {起点};
  \fill (N) circle (1.5pt) node[above=1pt] {终点};
  \draw (E) ++(0,0.35) -- ++(-0.35,0) -- ++(0,-0.35);
\end{tikzpicture}
\end{center}

向量家族在物理里的前四位成员，值得逐个点名。\textbf{位移}已经认识。\textbf{速度}是“位移变化的快慢加方向”：汽车仪表盘只显示快慢（速率，标量），导航还要告诉你朝哪开；第 5 章会正式定义它。\textbf{加速度}是“速度变化的快慢加方向”——注意，既然速度带着方向，那么\emph{光拐弯不变快慢也算加速}：方向盘打弯时你身体被甩向一边，就是这笔账在提醒你，这个伏笔会在圆周运动里长成主角。\textbf{力}在第 6 章已经登场，现在补上它的向量身份：推门，效果取决于用多大劲\emph{和}朝哪儿推，斜着推门的上沿，门可能又转又卡。

速度合成在生活中几乎随处可见，只是我们不叫它向量加法。在向上的自动扶梯上你也向上走，相对地面的快慢是“扶梯速度加步行速度”——同向，箭头直接相接；你要是在扶梯上向下走且走得和扶梯一样快，两支箭头反向相消，你就原地悬停在半空中，旁人会看到一幅诡异的画面：你的腿在走，人却不挪窝。飞机更是天天做这道算术：目的地在正北，风从正西吹来，机头如果只指正北，飞机会被风推着偏向东北；飞行员必须把机头偏向西北一个角，让“飞机相对空气”与“空气相对地面”两支箭头的合箭头恰好指向正北。你在航班屏幕上看到的航向和航迹常常不一致，那不是飞机开歪了，是向量加法在上班。

\step{一句话模型}

\begin{oneline}
带方向的量，按“首尾相接的箭头”做加减，方向本身就是数值的一部分；想知道它在某个指定方向上“有多少”，就把箭头往那个方向投影——每个方向各算各的账，互不干扰。
\end{oneline}

前半句管“怎么合成”，后半句管“怎么拆开”。合成与分解，就是本章全部的两件武器。

\step{数学翻译器}

\begin{mainline}
先把加法算熟。两支同样五十牛的箭头，首尾相接：同向时，合箭头长一百牛；反向时，恰好接回出发点，合箭头长度为零；夹成直角时，合箭头是直角三角形的斜边，长 \(\sqrt{50^{2}+50^{2}}\approx 71\) 牛。夹角越大，合箭头越短——这不是新物理，只是“三角形两边之和大于第三边”换了一身衣服。开场第三题的六十度角，答案八十七牛，一会儿在数学桥里用公式核对。

再看一眼拼路实验里藏着的勾股数：三步、四步、斜边五步，\(3^{2}+4^{2}=5^{2}\)，这是直角三角形最出名的一组整数解。你在家里量出的绳长如果正好是五步，那不是巧合，是勾股定理在客厅地板上值班。也请顺手做个“方向反转检查”：向北四步改成向南四步，合位移就从“东北方五步”变成“东南方五步”，大小没变、方向翻了——位移的方向反转，合向量跟着反转，这正是向量区别于标量的脾气。

再学分解。任何一支箭头都可以拆成沿两个垂直方向的两支，而且拆法唯一。为什么要拆？因为拆完之后，\emph{每个方向可以单独算账}。竖直方向的账不会改变水平方向的账，反过来也一样——这个“互不干扰”性质是整个经典力学最省心的礼物，第 5 章的抛体运动（水平匀速、竖直加速，各算各的）全靠它。

拉杆箱是分解的最佳教材。你的拉力 \(\vec{F}\) 沿拉杆斜向前上方，与水平面夹角记作 \(\theta\)。把它拆成两支：向前的分量大小是 \(F\cos\theta\)，向上的分量是 \(F\sin\theta\)。真正拖着箱子前进的只有向前那支；向上那支不推箱子，只负责把箱子往起抬一点、减轻它对地面的压力（顺带减小摩擦）。拉杆越陡，\(\theta\) 越大，\(\cos\theta\) 越小——同样的力气，向前的劲头越小。先猜一猜里的第三题，答案已经露头了。

分解也解释抬重物的配合问题。两个人一左一右抬一只沉柜子，如果两人的手臂都竖直向上用力，柜子重量由两人均摊，谁也不吃亏；可现实中两人站得开、手臂向中间倾斜，于是每个人使的劲都要拆成两支：向上的分量负责扛重量，水平向内或向外的分量互相顶着、纯属内耗。站得越开，手臂越斜，\(\cos\theta\) 越小，想凑够同样大的向上分量，每个人就得使更大的总劲——所以搬家公司的人抬东西总尽量贴近站。账还是那两笔账：每个方向各算各的。
\end{mainline}

\begin{center}
\begin{tikzpicture}[scale=1.0]
  \draw[thick,fill=gray!15] (0,0) rectangle (1.6,1.0);
  \node at (0.8,0.5) {箱子};
  \draw[thick] (-1.2,0) -- (5.4,0);
  \fill[pattern=north east lines] (-1.2,-0.25) rectangle (5.4,0);
  \coordinate (P) at (1.6,1.0);
  \draw[-{Stealth[length=3mm]},very thick,red!70!black] (P) -- ++(40:3.2) node[right=2pt] {拉力 \(\vec{F}\)};
  \draw[-{Stealth[length=3mm]},very thick,blue!70!black] (P) -- ++(0:2.45) node[midway,below=2pt] {向前分量 \(F\cos\theta\)};
  \draw[-{Stealth[length=3mm]},very thick,blue!70!black] (P) ++(0:2.45) -- ++(90:2.06);
  \node[right=2pt] at ($(P)+(2.45,1.2)$) {向上分量 \(F\sin\theta\)};
  \draw[dashed] (P) -- ++(0:1.5);
  \draw ($(P)+(0.9,0)$) arc (0:40:0.9) node[midway,right=1pt] {\(\theta\)};
\end{tikzpicture}
\end{center}

\begin{mainline}
“沿某个方向有多少”这件事太常用了，值得拥有姓名和公式。给两个向量 \(\vec{A}\) 和 \(\vec{B}\)，夹角为 \(\theta\)，定义
\[
  \vec{A}\cdot\vec{B} = AB\cos\theta ,
\]
读作“\(\vec{A}\) 的大小，乘以 \(\vec{B}\) 在 \(\vec{A}\) 方向上的投影”（等价地，\(\vec{B}\) 的大小乘以 \(\vec{A}\) 的投影，两种读法结果相同）。它叫\emph{点积}，结果是一个标量——方向的信息用完了就退场，留下纯粹的“多少”。影子实验就是点积的现场版：影长等于铅笔长度乘以 \(\cos\theta\)，\(\theta\) 是铅笔与光线的夹角。

新公式照例过三道检查。第一，\(\theta=0\)：\(\cos 0^{\circ}=1\)，点积等于 \(AB\)——完全同向时“有多少”就是全部，合理。第二，\(\theta=90^{\circ}\)：\(\cos 90^{\circ}=0\)，点积为零——竖直的铅笔在水平光里没有影子，垂直方向上“一点儿账都没有”，合理。第三，\(\theta=180^{\circ}\)：\(\cos 180^{\circ}=-1\)，点积为 \(-AB\)——负号不是印刷错误，它在说“方向拧反了”：你顶着风走，风在你前进方向上的账就是负的。

点积的第一个大用场在第 9 章：力推着物体前进一段位移，\emph{功}就是力与位移的点积——只有沿运动方向的那部分力才在“做功”的账本上签字。今天先记住它的样子：点积回答“一个方向上有多少”。
\end{mainline}

\begin{mathbridge}
有了坐标，箭头就变成了数对。选定两个垂直方向当坐标轴，向量 \(\vec{A}\) 写成 \((A_x,A_y)\)。加法是各分量分别相加：
\[
  \vec{A}+\vec{B} = (A_x+B_x,\;A_y+B_y),
\]
长度用勾股定理找回：\(A=\sqrt{A_x^2+A_y^2}\)。坐标轴可以随便转——向量本身不变，变的只是它的“读数”；选对方向，问题常常自己解体（比如沿斜面和垂直斜面分解重力，是第 7 章的常客）。

两个五十牛的力夹六十度角，用余弦定理核对开场的谜题：
\[
  F_{\text{合}}=\sqrt{50^{2}+50^{2}+2\times 50\times 50\cos 60^{\circ}}
  =\sqrt{7500}\approx 87\ \text{牛}.
\]
检查特例：夹角零度时根号里变成 \((50+50)^2\)，得一百牛；夹角一百八十度时变成 \((50-50)^2\)，得零。公式在三个极端上全都回归常识。

点积在坐标里更省事，连角度都不用知道：
\[
  \vec{A}\cdot\vec{B} = A_xB_x + A_yB_y .
\]
两种写法 \(AB\cos\theta\) 与 \(A_xB_x+A_yB_y\) 相等，用余弦定理展开 \(|\vec{A}-\vec{B}|^2\) 就能互相推出。这也回答了影子实验里你量到的影长变化：影长与铅笔长之比，正好描出 \(\cos\theta\) 的曲线。

最后看一眼一维：当所有向量都挤在同一条直线上，向量退化成“带正负号的数”。你小学的负数，其实是向量算术的幼年形态——方向只剩下“正”和“反”两种选择时的向量。
\end{mathbridge}

\begin{challenge}
点积问“一个方向上有多少”，还有另一种量问“这股劲有多想让你转起来”。开门的经验人人都有：力大不一定拧得动门——贴着门轴推，再大力也白搭；捏着门把手的远端、垂直于门板推，轻轻一碰门就开。决定转动效果的是三件事的乘积：力的大小 \(F\)、转轴到受力点的距离 \(r\)，以及方向因子 \(\sin\theta\)（\(\theta\) 是 \(\vec r\) 与 \(\vec F\) 的夹角）。物理学家把它打包成一个新的向量——\emph{叉积}
\[
  \vec{r}\times\vec{F},\qquad \text{大小 } rF\sin\theta ,
\]
方向用右手定则给出：四指从 \(\vec r\) 弯向 \(\vec F\)，拇指指向结果。奇怪但深刻的是，这个“转动倾向”的量不在 \(\vec r\) 和 \(\vec F\) 张成的平面里，而是垂直地戳出来——这只有在三维空间才办得到，也是它直到三维向量算术成熟后才被发明的原因。照例检查特例：\(\theta=0\)（沿扳手方向推拉）结果为零，拧不动，对；\(\theta=90^{\circ}\) 取到最大值 \(rF\)，垂直发力最有效，对。叉积是第 11 章的主角力矩和角动量的数学骨架，那里它还要负责解释陀螺为什么晃而不倒、花滑选手收臂为什么越转越快。

还有一件事值得在此预告：向量不一定住在空间里。凡是“几个数排成一队、按分量加减、能整体放大缩小”的对象，都满足向量的全部算术规则。到第 46 章你会看到，量子力学把一个系统的整个状态当成一支向量——只是那支箭头不再指向东南西北，而是指向抽象状态空间里的某个方向。本章学会的“分解到选定方向上各算各的账”，到那时会原封不动地再用一次。
\end{challenge}

\step{模型的边界}

\begin{mainline}
向量算术在日常生活尺度上好用到让人忘记它有边界。以下是它的国界和本章特有的误用清单。

边界一：\emph{速度合成在近光速时失效}。“船速加水速”式的加法，在速度远小于光速时精确得无可挑剔；但当速度接近光速，实验发现 \(0.8\) 倍光速叠加 \(0.8\) 倍光速不等于 \(1.6\) 倍光速——任何叠加都越不过光速。这不是向量算术错了，而是高速世界里时间和空间本身的关系要重写，加法公式需要修正，详见第 38、39 章。日常速度下修正小到测不出，所以本章照常使用。

边界二：\emph{有限的转动不是向量}。拿一本书做个十秒实验：先绕水平轴向前翻九十度，再绕竖直轴向左转九十度，记住书的姿态；然后从原姿态重来，把两步次序对调。两次的结果不一样！带方向的量按向量加法合成时，次序无关紧要（\(\vec A+\vec B=\vec B+\vec A\)），而有限转动的结果依赖次序，所以“转九十度”不能当成一支箭头。注意区分：转动的\emph{快慢}（角速度）是向量，有限的\emph{转角}不是。第 11 章会收拾这个烂摊子。

边界三：\emph{球面上的方向}。本章的箭头默认活在平面上。在地球这样的球面上走长距离，“一直向东”会沿纬线绕回原地，而两点之间最直的路线（大圆航线）在地图上是弯的——远程航班的航线因此总要“向北绕”。平面向量是局部近似，够我们用在城市、操场和实验室里。

边界零（之所以排最前，因为它最常被违反）：\emph{只有同一种量才能相加}。位移加位移、力加力，都有意义；“三米加五牛”不是算不出来，而是根本没有这支箭头可画。向量加法不改变第 2 章的规矩：先把量纲对上台账，再谈方向怎么拼。反过来，这个约束也是免费的查错器——如果一道题的中间步骤里出现了“速度加上力”，不用往下算，前面一定写错了。
\end{mainline}

典型误用，每条都对应前面埋过的雷：

\begin{itemize}[leftmargin=1.8em]
  \item \emph{把大小相加当成向量相加。}“他用了五十牛，我也用了五十牛，所以总共一百牛”——先问方向。大小直接相加只在完全同向时成立。
  \item \emph{把路程当位移，或反过来。}罚跑操场一圈，路程四百米，位移是零；导航说“距目的地五百米”说的是位移大小，不是你还要走的路程。
  \item \emph{以为分量“不是真的”。}向前的分量和向上的分量不是两个假想力，而是同一个拉力在两个方向上的真实账目；但反过来，把一个力拆成分量也不等于“多出来两个力”——分解是记账方式，不是加戏。
  \item \emph{以为向量大的效果一定大。}斜着拉杆箱，力气大不如角度好。效果取决于你要的是哪个方向上的账。
\end{itemize}

\step{回头解释开场现象}

\textbf{雨痕。}新模型的回答是：车窗上水痕的方向，由雨\emph{相对于车}的速度决定，而相对速度是向量减法——雨相对车的速度，等于雨相对地面的速度，减去车相对地面的速度。雨竖直向下，车水平向前；从车的视角看，雨不但向下落，还迎面向后“飞”来。两个方向首尾相接一减，合出来的箭头斜向前下方，于是雨斜着砸向侧窗，水痕从前上方拖向后下方——你在先猜一猜里如果画了“上端靠车头、下端靠车尾”，猜对了。

来一笔带真实数据的估算。大雨雨滴下落的收尾速度约每秒八米（毛毛雨约每秒两米，暴雨的大雨滴可达每秒九米），汽车以每小时六十公里行驶，约每秒十七米。水痕与竖直方向的夹角满足
\[
  \tan\theta=\frac{v_{\text{车}}}{v_{\text{雨}}}=\frac{17}{8}\approx 2.1,
  \qquad \theta\approx 65^{\circ} .
\]
也就是说，时速六十公里时雨痕与竖直方向偏了约六十五度，难怪看起来“快横过来了”。检查特例：车停下，\(v_{\text{车}}=0\)，\(\theta=0\)，雨痕竖直，对；车速趋于极大，\(\theta\) 趋近九十度，雨痕趋近水平，对。从头到尾没有谁“掰”雨——斜的不是雨，是你和雨之间的相对运动。

\textbf{拼路。}七百米是路程，标量加法；五百米是位移大小，向量加法。“消失的两百米”是个伪问题：两笔账规则不同，答案当然不同。如果你在原地绕圈走了一千米，计步器说一千米，位移是零——它一米都没“消失”，因为位移从来不承诺汇报你出过多少力。

\textbf{两只弹簧秤。}五十牛加五十牛，夹角六十度，合拉力按向量合成约八十七牛，数学桥里已经用余弦定理核对。弹簧秤没有违抗加法，违抗的是“力是标量”的旧假设。力是向量，它服从的是箭头加法——这是用无数实验投票通过的实验事实。

三个悬念就此销账，顺手回顾一下你在家里做的两个实验：拼路实验就是位移加法本身，绳长五步是勾股定理；影子实验就是点积的雏形，影长与铅笔长之比随角度画出的正是 \(\cos\theta\) 的曲线。你已经亲手摸过了本章的全部两件武器。

新模型还顺手送了一件赠品：解释你没问过的问题。\textbf{帆船为什么能逆风前进，甚至跑得比风快？}帆面斜对来风时，风推帆的力大致垂直于帆面；把这个力分解到船头方向和侧向，向前的那支推船前进，侧向的那支被船底的龙骨或稳向板顶住。于是正对风来的方向反而是帆船的禁区（帆只抖动不受力，水手叫它“无法航行区”），聪明的水手走之字形路线，斜着迎风一程、换舷再一程，像爬楼梯一样折向风来的方向挤上去——每一段的合速度都是船速与风压方向的向量游戏。更惊人的是冰船和陆地帆船：阻力极小时，船的行进速度可以明显超过风速本身，二〇〇九年陆地帆船“绿鸟”号在美国内华达州的干涸湖床上跑出约每小时二百零三公里，是当时风速的三倍上下。直觉说“风推的船怎么可能比风快”，向量说：谁也没规定船的速度必须是风速的一个分量。只要阻力足够小、帆的角度调得对，船自己的速度反过来会改变它“感受到”的风向，让帆继续吃到向前的分量——这笔正反馈的账，留给你在挑战题之外慢慢算。

\step{脑洞实验与挑战题}

\begin{thought}[你此刻正在以每秒二百三十公里运动]
把向量合成推到极端尺度。你坐在椅子上觉得自己纹丝不动，可“速度”是相对于参考系的向量：地球赤道上的自转带着你以约每秒零点五公里向东；地球绕太阳公转带着你以约每秒三十公里沿轨道切线飞；太阳又带着整个太阳系绕银河系中心以约每秒二百三十公里狂奔。你此刻“真实”的速度，是这三支箭头的向量和——而且因为前两个方向天天在变，这个和的大小和方向每时每刻都在改写：午夜和正午，你相对太阳系的朝向相差将近一百八十度。“静止”从来不是一个绝对状态，它只是一句省略了参考系的口头语。下次有人说“坐着不动最稳”，你可以告诉他：不动的是椅子相对于地板的那一支向量，其余的账，宇宙还在哗哗地记。
\end{thought}

\begin{thought}[星光是一场雨：把车窗搬到天文台]
车窗雨痕的思路放大到天文学，会得出一个真实的科学发现。遥远恒星的光垂直“落”向太阳系，地球却以每秒约三十公里绕着太阳跑——地球就是那辆车，星光就是那场雨。于是望远镜里的星光方向应当像雨痕一样倾斜：星光相对地球的速度，是光速向量减去地球公转速度向量。这个效应叫\emph{光行差}，一七二九年布拉德雷真的测到了它——所有恒星每年在天上画一个小椭圆，幅度与“地球速度比光速”这个比值精确吻合，他还顺手用它算出了光速。但故事有个漂亮的尾巴：当测量精度继续提高，经典向量加法预言的偏角与实际出现了微小而系统的偏差。修这个偏差不能用更粗的箭头，而需要重写速度合成规则本身——那是第 38、39 章相对论的入口。本章的一支小箭头，一路指向了二十世纪的物理学革命。
\end{thought}

\begin{puzzles}
\item 河水以每秒一米的速度向下游流，小船相对水的速度为每秒两米。船夫把船头正对对岸（船相对水的方向与河岸垂直）划，船会在正对岸登陆吗？若想在正对岸登陆，船头该朝哪个方向偏？\emph{思路提示}：船相对岸的速度是“船相对水”与“水相对岸”两支箭头的和。正对岸划时，合箭头斜向下游，登陆点在下游。想正对岸登陆，得让合箭头垂直河岸——船头要朝上游偏，让船速向上游的分量恰好抵消水流。两支分量，各算各的账。
\item 保持拉力大小不变，拉杆箱的把手抬得越陡，向前的分量怎样变化？可另一方面，抬陡一点箱子对地面的压力变小、摩擦也变小——为什么仍然存在一个“最舒服的角度”，而不是越陡越好或越平越好？\emph{思路提示}：向前分量是 \(F\cos\theta\)，随角度变陡单调变小；摩擦的负担随 \(\sin\theta\) 增大而减轻。两笔账一笔减一笔增，舒服的角度是两笔账的平衡点。工程上行李箱拉杆的角度设计，做的就是这笔向量账。
\item 两个人各用一百牛的力、夹角一百二十度拉同一个小环。现在来第三个人，他至少要用多大的力、朝什么方向拉，才能让环纹丝不动？\emph{思路提示}：先求前两支箭头的合向量。夹角一百二十度有个巧合：由余弦定理，合力大小恰好也是一百牛，方向沿两力的角平分线。第三个人只需与这个合力等大反向。注意三个人用的力一样大——对称不是偶然，是几何。
\item 本章说“有限的转动不是向量”，证据是两次九十度翻转的次序不能交换。可是角速度却被当成向量（比如地球自转轴指向北极星附近）。同一个“转动”，为什么一个行一个不行？\emph{思路提示}：关键在“无限小”。把每次翻转切成无数次无限小的转动，次序造成的影响会趋于零——无限小转动的合成不依赖次序，可以用箭头表示；而任何有限大的转动都残留次序的痕迹。这个“无限小才守规矩”的性质，正是第 3 章里极限思想在转动问题上的重演。
\end{puzzles}
